%==========================================================================
% Segunda Lista de Engenharia de Software RT
% Capítulo 2 — Ian Sommerville (trabalho individual)
%==========================================================================
\documentclass[12pt,a4paper]{article}

\usepackage[T1]{fontenc}
\usepackage[utf8]{inputenc}
\usepackage[brazil]{babel}
\usepackage{amsmath}
\usepackage{geometry}
\usepackage{xcolor}
\usepackage{tcolorbox}
\usepackage{enumitem}
\usepackage{tikz}
\usetikzlibrary{shapes.geometric,arrows.meta,positioning,calc}
\usepackage{hyperref}
\usepackage{booktabs}
\usepackage{array}

\geometry{top=2.5cm, bottom=2.5cm, left=3cm, right=3cm}
\setlength{\parskip}{6pt}
\setlength{\parindent}{0pt}

\tcbuselibrary{skins,breakable}
\newtcolorbox{questao}[1]{colback=gray!8,colframe=gray!50,
  fonttitle=\bfseries,title=#1,breakable}
\newtcolorbox{resposta}{colback=blue!4,colframe=blue!50!black,
  fonttitle=\bfseries,title=Resposta,breakable}
\newtcolorbox{nota}{colback=orange!5,colframe=orange!60!black,
  fonttitle=\bfseries,title=Nota,breakable,before upper=\small}

%==========================================================================
\begin{document}

\begin{center}
  {\Large\bfseries Segunda Lista de Engenharia de Software RT}\\[6pt]
  {\large Capítulo 2 — Processos de Software}\\[4pt]
  {\normalsize Ian Sommerville — \textit{Software Engineering}, 10ª ed.}\\[2pt]
  {\small Trabalho Individual}
\end{center}
\hrule\vspace{8pt}

%==========================================================================
\section*{Questão 1}
%==========================================================================

\begin{questao}{Enunciado}
Justificando sua resposta com base no tipo de sistema a ser desenvolvido,
sugira o modelo genérico de processo de software mais adequado para ser
usado como base para a gerência do desenvolvimento dos sistemas a seguir:
\begin{enumerate}[label=\alph*.]
  \item Um sistema para controlar o antibloqueio de frenagem de um carro (ABS).
  \item Um sistema de realidade virtual para dar apoio à manutenção de software.
  \item Um sistema de contabilidade para uma universidade, que substitua um sistema já existente.
  \item Um sistema interativo de planejamento de viagens que ajude os usuários a
        planejar viagens com menor impacto ambiental.
\end{enumerate}
\end{questao}

\begin{resposta}

\textbf{(a) Sistema ABS — Modelo em Cascata (Waterfall)}

Um sistema de antibloqueio de frenagem é um \textbf{sistema crítico de tempo real}
embarcado com requisitos de segurança altamente rigorosos e definidos antes do
início do desenvolvimento. Falhas podem custar vidas. Por isso, exige:
\begin{itemize}[nosep]
  \item Especificação completa e formal de requisitos antes da implementação;
  \item Validação e verificação formais em cada fase;
  \item Rastreabilidade total entre requisitos, projeto, código e testes.
\end{itemize}

O \textbf{modelo em cascata} é o mais adequado porque impõe uma sequência rígida
de fases (requisitos → projeto → implementação → verificação → manutenção) com
saídas documentadas em cada etapa, atendendo a padrões como IEC~61508 (segurança
funcional). A completeza dos requisitos antes da codificação é viável pois os
requisitos físicos do sistema são bem compreendidos e estáveis.

\vspace{8pt}
\textbf{(b) Sistema de Realidade Virtual para Manutenção de Software — Desenvolvimento Evolucionário / Espiral}

Sistemas de realidade virtual são inerentemente exploratórios: as interfaces,
metáforas de interação e as necessidades dos usuários só ficam claras na prática.
Os requisitos \emph{não podem ser completamente especificados} antecipadamente.

O \textbf{modelo evolucionário} (ou o \textbf{modelo espiral de Boehm}) é mais
adequado porque:
\begin{itemize}[nosep]
  \item Permite protótipos sucessivos com avaliação de usuários a cada ciclo;
  \item Incorpora gerência de riscos explícita (fundamental em tecnologia nova);
  \item Evoluções do sistema são esperadas conforme os usuários descobrem novas necessidades.
\end{itemize}

\vspace{8pt}
\textbf{(c) Sistema de Contabilidade Universitário (substituição) — Desenvolvimento Baseado em Reuso / Incremental}

Como já existe um sistema legado em operação, os requisitos são relativamente
conhecidos (mesmo fluxo contábil) e há grande disponibilidade de componentes
reutilizáveis (frameworks ERP, módulos de relatórios fiscais, bibliotecas
contábeis). O risco principal é a migração de dados e a aceitação dos usuários.

O \textbf{modelo baseado em reuso} é o mais adequado porque:
\begin{itemize}[nosep]
  \item Maximiza o uso de software pronto (reduz custo e prazo);
  \item Entregas incrementais permitem migração gradual com menor risco operacional;
  \item Requisitos derivam do sistema existente, minimizando incertezas.
\end{itemize}

\vspace{8pt}
\textbf{(d) Sistema Interativo de Planejamento de Viagens — Desenvolvimento Incremental / Ágil}

Este é um sistema voltado a usuários finais em um domínio em constante evolução
(regulamentações ambientais, bases de dados de emissões, preferências de viajantes).
Os requisitos mudarão ao longo do desenvolvimento conforme os usuários interagem
com o produto.

O \textbf{desenvolvimento incremental} (base para métodos ágeis) é o mais adequado:
\begin{itemize}[nosep]
  \item Entregas frequentes de funcionalidades priorizadas;
  \item Feedback contínuo dos usuários incorporado em cada incremento;
  \item Facilidade de adaptação a novas fontes de dados ambientais ou regulamentações.
\end{itemize}

\end{resposta}

%==========================================================================
\section*{Questão 2}
%==========================================================================

\begin{questao}{Enunciado}
Explique por que o desenvolvimento incremental é o método mais eficaz para
o desenvolvimento de sistemas de software de negócios. Por que esse modelo
é menos adequado para a engenharia de sistemas de tempo real?
\end{questao}

\begin{resposta}

\textbf{Eficácia no desenvolvimento de sistemas de negócios:}

Sistemas de negócios operam em ambientes dinâmicos onde requisitos mudam
continuamente em resposta a alterações regulatórias, de mercado e
organizacionais. O desenvolvimento incremental é mais eficaz porque:

\begin{enumerate}[nosep]
  \item \textbf{Entrega de valor antecipada:} versões funcionais do sistema são
        disponibilizadas cedo, gerando retorno sobre o investimento antes da
        conclusão total do projeto.
  \item \textbf{Validação contínua:} usuários de negócio verificam cada incremento
        e fornecem feedback, reduzindo o risco de o sistema final não atender
        às necessidades reais.
  \item \textbf{Acomodação de mudanças:} como os requisitos de negócio evoluem,
        incrementos futuros absorvem as mudanças sem comprometer o que já foi entregue.
  \item \textbf{Redução do risco financeiro:} problemas são detectados cedo, quando
        o custo de correção ainda é baixo.
  \item \textbf{Gerenciamento do escopo:} funcionalidades prioritárias são entregues
        primeiro, garantindo que o sistema seja útil mesmo se o projeto não for
        concluído integralmente.
\end{enumerate}

\textbf{Inadequação para sistemas de tempo real:}

Sistemas de tempo real (controle industrial, aviônicos, sistemas médicos)
têm \textbf{requisitos não-funcionais críticos} --- latência, taxa de resposta,
determinismo temporal --- que não podem ser validados em incrementos isolados:

\begin{enumerate}[nosep]
  \item \textbf{Integração indivisível:} o comportamento temporal emergente só pode
        ser verificado quando todos os componentes estão integrados. Um incremento
        parcial pode satisfazer requisitos funcionais mas violar garantias de tempo.
  \item \textbf{Certificação formal:} normas de segurança (DO-178C, IEC~61508)
        exigem documentação e rastreabilidade completas desde os requisitos, o que
        conflita com a natureza evolutiva do modelo incremental.
  \item \textbf{Hardware acoplado:} o comportamento de tempo real depende da
        interação precisa entre software e hardware; mudanças incrementais no
        software podem invalidar análises de \textit{worst-case execution time}.
  \item \textbf{Custo de falha catastrófico:} em sistemas críticos, uma falha
        entregue em produção pode custar vidas, inviabilizando a entrega precoce
        de versões incompletas.
\end{enumerate}

\end{resposta}

%==========================================================================
\section*{Questão 3}
%==========================================================================

\begin{questao}{Enunciado}
Considere o modelo de processo baseado em reuso da Figura~2.3. Explique por
que, nesse processo, é essencial ter duas atividades distintas de engenharia
de requisitos.
\end{questao}

\begin{resposta}

No modelo baseado em reuso, o processo possui duas atividades de engenharia
de requisitos:

\begin{enumerate}
  \item \textbf{Especificação de Requisitos Iniciais} — captura o que o cliente
        \emph{realmente precisa}, sem considerar o que está disponível para reuso.
        É necessária para estabelecer uma linha de base das necessidades do sistema.

  \item \textbf{Modificação de Requisitos} — após a descoberta e avaliação dos
        componentes reutilizáveis disponíveis, os requisitos originais são
        \emph{revisados e adaptados} para refletir o que os componentes
        conseguem oferecer.
\end{enumerate}

\textbf{Por que as duas atividades são essenciais:}

A primeira especificação (independente de reuso) é necessária porque os
componentes disponíveis são descobertos \emph{após} a elicitação inicial. Sem
ela, não é possível avaliar quais requisitos podem ser satisfeitos por reuso e
quais precisam de desenvolvimento novo.

A segunda atividade existe porque os componentes reutilizados raramente satisfazem
os requisitos originais \textbf{exatamente}. Há um \textit{gap} entre o que o
cliente deseja e o que o componente oferece. Ajustar os requisitos ao invés de
forçar a adaptação do componente preserva a integridade arquitetural e reduz
custos de manutenção. Sem essa segunda fase, o sistema seria desenvolvido com
discrepâncias não documentadas entre requisitos e funcionalidade real, gerando
insatisfação do cliente e dívida técnica.

\begin{nota}
Sommerville destaca que a principal desvantagem do modelo baseado em reuso é
exatamente essa negociação de requisitos: o cliente pode precisar aceitar um
sistema que não atende completamente suas necessidades originais.
\end{nota}

\end{resposta}

%==========================================================================
\section*{Questão 4}
%==========================================================================

\begin{questao}{Enunciado}
Sugira porque é importante, no processo de engenharia de requisitos, fazer
uma distinção entre desenvolvimento dos requisitos do usuário e
desenvolvimento de requisitos de sistema.
\end{questao}

\begin{resposta}

\textbf{Requisitos do usuário} descrevem, em linguagem natural acessível, o que
os usuários e demais \textit{stakeholders} esperam que o sistema faça para apoiar
seu trabalho. São \textbf{orientados a objetivos} e independentes de como o sistema
será implementado.

\textbf{Requisitos de sistema} são especificações técnicas detalhadas, precisas e
completas de \emph{como} o sistema deve realizar o que foi pedido, incluindo
restrições de desempenho, protocolos, interfaces, segurança e consistência de dados.

\textbf{Por que a distinção é importante:}

\begin{enumerate}
  \item \textbf{Audiências diferentes:} gerentes, clientes e usuários finais leem
        os requisitos de usuário para validar que o sistema resolve seus problemas;
        arquitetos e desenvolvedores leem os requisitos de sistema para implementá-lo.
        Misturar os dois níveis confunde cada grupo.

  \item \textbf{Comprometimento contratual:} os requisitos de usuário costumam
        compor o contrato entre cliente e fornecedor. Requisitos de sistema são
        internos à equipe de desenvolvimento. Confundir os dois pode gerar disputas
        sobre o que foi contratado.

  \item \textbf{Rastreabilidade:} a separação permite rastrear cada requisito de
        sistema a um ou mais requisitos de usuário, facilitando a verificação de
        que todos os objetivos do cliente foram atendidos.

  \item \textbf{Gerenciamento de mudanças:} mudanças em requisitos de usuário têm
        impacto diferente de mudanças em requisitos de sistema. Separar os dois
        permite avaliar o impacto de cada mudança com precisão.

  \item \textbf{Detalhamento progressivo:} permite um processo de engenharia
        gradual — primeiro valida-se o \emph{o quê} com os stakeholders, depois
        especifica-se o \emph{como} com a equipe técnica, evitando refazer trabalho
        técnico por mal entendimento de objetivos.
\end{enumerate}

\end{resposta}

%==========================================================================
\section*{Questão 5}
%==========================================================================

\begin{questao}{Enunciado}
Descreva as principais atividades do processo de projeto de software e as
saídas dessas atividades. Usando um diagrama, mostre as possíveis relações
entre as saídas dessas atividades.
\end{questao}

\begin{resposta}

As principais atividades do processo de projeto de software, conforme
Sommerville, são:

\begin{center}
\begin{tabular}{p{4.5cm}p{8.5cm}}
\toprule
\textbf{Atividade} & \textbf{Saída} \\
\midrule
Projeto de arquitetura & Diagrama de componentes/módulos; estilo arquitetural
                         (camadas, cliente-servidor, MVC etc.) \\
\midrule
Projeto de banco de dados & Modelo entidade-relacionamento; esquema lógico/físico;
                            dicionário de dados \\
\midrule
Projeto de interface & Protótipos de telas; descrição de fluxo de navegação;
                       especificação de API/protocolo de comunicação \\
\midrule
Projeto de componentes & Especificação detalhada de cada módulo; algoritmos;
                         pré/pós-condições; pseudocódigo ou UML de classe \\
\bottomrule
\end{tabular}
\end{center}

\textbf{Diagrama de relações entre as saídas:}

\begin{center}
\begin{tikzpicture}[
    node distance=1.6cm and 2.4cm,
    box/.style={draw,rounded corners,fill=blue!8,minimum width=3.2cm,
                minimum height=1cm,align=center,font=\small},
    doc/.style={draw,fill=orange!10,minimum width=2.8cm,
                minimum height=0.8cm,align=center,font=\footnotesize},
    arr/.style={-{Stealth},thick},
    darr/.style={dashed,-{Stealth},gray}]

  % Activities
  \node[box] (arq) {Projeto de\\Arquitetura};
  \node[box, right=2.5cm of arq] (bd)  {Projeto de\\Banco de Dados};
  \node[box, below=1.6cm of arq] (iface) {Projeto de\\Interface};
  \node[box, below=1.6cm of bd] (comp)  {Projeto de\\Componentes};

  % Outputs
  \node[doc, below=0.5cm of arq]  (d_arq)  {Diagrama de\\Arquitetura};
  \node[doc, below=0.5cm of bd]   (d_bd)   {Esquema de BD\\+ Dicionário};
  \node[doc, below=0.5cm of iface](d_iface){Protótipos\\+ API spec};
  \node[doc, below=0.5cm of comp] (d_comp) {Especificação\\de Módulos};

  % Arrows activity -> output
  \draw[arr] (arq)   -- (d_arq);
  \draw[arr] (bd)    -- (d_bd);
  \draw[arr] (iface) -- (d_iface);
  \draw[arr] (comp)  -- (d_comp);

  % Relations between outputs
  \draw[darr] (d_arq)  -- node[above,font=\tiny]{refina} (d_bd);
  \draw[darr] (d_arq)  -- node[left,font=\tiny]{guia} (d_iface);
  \draw[darr] (d_bd)   -- node[right,font=\tiny]{detalha} (d_comp);
  \draw[darr] (d_iface)-- node[below,font=\tiny]{especifica} (d_comp);
  \draw[darr] (d_arq.south) .. controls +(0,-0.6) and +(0,0.6) ..
              node[right,font=\tiny,pos=0.5]{decompõe} (d_comp.north);
\end{tikzpicture}
\end{center}

O \textbf{diagrama de arquitetura} é a saída central: define os subsistemas e
suas interfaces, orientando todas as demais atividades. O projeto de banco de
dados refina a arquitetura de dados; o projeto de interface especifica os
contratos entre componentes; e o projeto de componentes detalha cada módulo
individualmente. Todas as saídas se integram na documentação final de projeto
que orienta a fase de implementação.

\end{resposta}

%==========================================================================
\section*{Questão 6}
%==========================================================================

\begin{questao}{Enunciado}
Explique por que, em sistemas complexos, as mudanças são inevitáveis.
Exemplifique as atividades de processo de software que ajudam a prever as
mudanças e fazer com que o software seja desenvolvido mais tolerante a
mudanças (desconsidere prototipação e entrega incremental).
\end{questao}

\begin{resposta}

\textbf{Por que mudanças são inevitáveis em sistemas complexos:}

\begin{enumerate}[nosep]
  \item \textbf{Ambiente externo mutável:} legislação, regulamentações,
        condições de mercado e tecnologias de suporte mudam independentemente
        da vontade do time de desenvolvimento.
  \item \textbf{Requisitos mal compreendidos:} em sistemas complexos, os
        stakeholders frequentemente só descobrem o que realmente precisam
        após ver o sistema em operação.
  \item \textbf{Conflito de interesses:} diferentes stakeholders têm visões
        conflitantes que só emergem durante o desenvolvimento.
  \item \textbf{Evolução do negócio:} a organização que solicitou o sistema pode
        se reestruturar, mudar de estratégia ou ser adquirida durante o projeto.
  \item \textbf{Descoberta de falhas:} erros nos requisitos e no projeto só
        são revelados em fases posteriores, exigindo revisões retroativas.
\end{enumerate}

\textbf{Atividades que ajudam a prever e tolerar mudanças}
(excluindo prototipação e entrega incremental):

\begin{itemize}
  \item \textbf{Antecipação de mudanças (prevenção):}
  \begin{itemize}[nosep]
    \item \textit{Revisões de requisitos (\textit{requirements reviews}):} análise
          sistemática dos requisitos com os stakeholders para identificar
          inconsistências, ambiguidades e itens provavelmente instáveis, antes
          da implementação.
    \item \textit{Rastreabilidade de requisitos:} manter matrizes de rastreabilidade
          entre requisitos e artefatos de projeto permite avaliar o impacto de
          mudanças antes de implementá-las.
  \end{itemize}

  \item \textbf{Tolerância a mudanças (acomodação):}
  \begin{itemize}[nosep]
    \item \textit{Projeto orientado a mudanças (\textit{change-oriented design}):}
          uso de princípios como alta coesão, baixo acoplamento e ocultamento de
          informação (\textit{information hiding}) isola módulos, de modo que uma
          mudança em um componente não propague efeitos para outros.
    \item \textit{Gerenciamento de configuração:} controle de versões, baselines e
          gestão de mudanças formais garantem que alterações sejam rastreadas e
          revertidas se necessário, sem comprometer a integridade do sistema.
  \end{itemize}
\end{itemize}

\end{resposta}

%==========================================================================
\section*{Questão 7}
%==========================================================================

\begin{questao}{Enunciado}
Explique por que os sistemas desenvolvidos como protótipos normalmente não
devem ser usados como sistemas de produção.
\end{questao}

\begin{resposta}

Protótipos são construídos sob restrições fundamentalmente diferentes das de um
sistema de produção:

\begin{enumerate}
  \item \textbf{Requisitos não-funcionais ignorados:} protótipos priorizam
        demonstrar funcionalidades visíveis. Desempenho, segurança, disponibilidade,
        escalabilidade e tolerância a falhas são deliberadamente omitidos para
        acelerar o desenvolvimento.

  \item \textbf{Dívida técnica intencional:} atalhos de implementação, ausência de
        tratamento de erros, hardcoding de valores e falta de modularização são
        aceitáveis em protótipos mas comprometem a manutenibilidade e confiabilidade
        em produção.

  \item \textbf{Ausência de testes rigorosos:} um protótipo raramente passa por
        ciclos completos de testes de sistema, regressão e segurança. Usá-lo
        em produção expõe os usuários a falhas não descobertas.

  \item \textbf{Tecnologias inadequadas:} linguagens, frameworks ou bancos de
        dados escolhidos pela velocidade de prototipagem podem não ser adequados
        para escala, manutenção de longo prazo ou integração com sistemas legados.

  \item \textbf{Ausência de documentação:} protótipos carecem de documentação de
        arquitetura, requisitos rastreáveis e documentação de manutenção,
        tornando a evolução futura extremamente custosa.

  \item \textbf{Pressão organizacional indevida:} transformar um protótipo em
        sistema de produção cria uma \textit{proto-trap} — a organização fica
        presa a decisões técnicas ruins tomadas sob pressão de prazo, pagando
        juros de dívida técnica indefinidamente.
\end{enumerate}

A conclusão de Sommerville é que protótipos devem ser \textbf{descartados} após
cumprir seu propósito de elicitação/validação de requisitos, e o sistema real
deve ser desenvolvido do zero com base nos aprendizados obtidos.

\end{resposta}

%==========================================================================
\section*{Questão 8}
%==========================================================================

\begin{questao}{Enunciado}
Explique por que o modelo em espiral de Boehm é um modelo adaptável que
apoia tanto as atividades de prevenção de mudanças quanto as de tolerância a
mudanças. Na prática, esse modelo não tem sido amplamente usado. Sugira as
possíveis razões para isso.
\end{questao}

\begin{resposta}

\textbf{O modelo espiral como modelo adaptável:}

O modelo espiral de Boehm organiza o processo em \textbf{ciclos iterativos},
cada um com quatro quadrantes: (1) definição de objetivos e alternativas,
(2) \emph{avaliação de riscos e prototipação}, (3) desenvolvimento e validação,
(4) planejamento da próxima iteração.

\begin{itemize}
  \item \textbf{Prevenção de mudanças:} a análise explícita de riscos em cada
        ciclo identifica antecipadamente áreas de incerteza (técnica, de
        requisitos, organizacional). Prototipação e avaliação antecipam problemas
        antes de grandes investimentos, funcionando como atividade de
        ``antecipação de mudanças''.

  \item \textbf{Tolerância a mudanças:} a estrutura iterativa permite que cada
        ciclo incorpore mudanças descobertas no ciclo anterior. Não há
        comprometimento irreversível com uma arquitetura específica no início
        do projeto; o design evolui com o crescente entendimento do problema.
\end{itemize}

\textbf{Por que o modelo espiral não é amplamente adotado na prática:}

\begin{enumerate}
  \item \textbf{Complexidade de gestão:} cada ciclo exige análise de riscos
        formal por profissionais especializados. Poucos times têm expertise
        em gerenciamento de riscos quantitativo, e a documentação necessária
        é extensa.

  \item \textbf{Custo e overhead:} a cerimônia de cada ciclo (documentação,
        revisão de riscos, planejamento) é custosa para projetos pequenos ou
        médios, nos quais a relação custo-benefício não se justifica.

  \item \textbf{Dificuldade contratual:} contratos de desenvolvimento de software
        geralmente exigem escopo, prazo e custo fixos. O modelo espiral pressupõe
        escopo emergente, o que é incompatível com contratos de preço fixo.

  \item \textbf{Surgimento de métodos ágeis:} a partir dos anos 2000, métodos
        ágeis (Scrum, XP, Kanban) ofereceram adaptabilidade similar com muito
        menos overhead, tornando o modelo espiral menos atrativo.

  \item \textbf{Falta de ferramentas de suporte:} não há ferramentas amplamente
        adotadas que automatizem ou guiem a avaliação de riscos periódica
        exigida pelo modelo espiral.
\end{enumerate}

\end{resposta}

%==========================================================================
\section*{Questão 9}
%==========================================================================

\begin{questao}{Enunciado}
Quais são as vantagens de proporcionar visões estáticas e dinâmicas do
processo de software, assim como no Rational Unified Process?
\end{questao}

\begin{resposta}

O \textbf{Rational Unified Process (RUP)} oferece simultaneamente:

\begin{itemize}
  \item \textbf{Visão dinâmica (eixo temporal):} organizada em quatro fases ---
        Iniciação, Elaboração, Construção e Transição --- que descrevem
        \emph{quando} as atividades ocorrem ao longo do projeto, com marcos
        e objetivos específicos em cada fase.

  \item \textbf{Visão estática (fluxos de trabalho):} organizada em disciplinas
        --- Requisitos, Análise e Projeto, Implementação, Testes, Implantação
        etc. --- que descrevem \emph{o que} é feito, independentemente do tempo.
\end{itemize}

\textbf{Vantagens da combinação de ambas as visões:}

\begin{enumerate}
  \item \textbf{Planejamento e execução alinhados:} a visão dinâmica permite ao
        gerente de projeto planejar cronogramas e marcos. A visão estática garante
        que os engenheiros saibam quais práticas aplicar em cada disciplina.

  \item \textbf{Comunicação clara com diferentes audiências:} gestores entendem
        o processo pelo eixo temporal (fases, entregas); engenheiros o entendem
        pelo eixo das disciplinas (o que fazer e como). As duas visões atendem
        a ambos sem perda de informação.

  \item \textbf{Adaptabilidade controlada:} todas as disciplinas ocorrem em
        todas as fases, mas com intensidade variável. Isso deixa explícito que
        testes, por exemplo, não são atividades apenas do final do projeto, sem
        perder a estrutura temporal que orienta prioridades.

  \item \textbf{Rastreabilidade:} a intersecção fase $\times$ disciplina
        permite identificar com precisão quando e onde artefatos específicos
        devem ser produzidos, facilitando auditorias e revisões.

  \item \textbf{Suporte a iteratividade:} as fases dinâmicas acomodam iterações
        dentro de cada fase, enquanto as disciplinas estáticas garantem que
        boas práticas de engenharia sejam aplicadas consistentemente em cada
        iteração, independentemente do ponto do projeto.
\end{enumerate}

\end{resposta}

%==========================================================================
\section*{Questão 10}
%==========================================================================

\begin{questao}{Enunciado}
Historicamente, a introdução de tecnologia provocou mudanças profundas no
mercado de trabalho e, pelo menos temporariamente, deixou muitas pessoas
desempregadas. Discuta se a introdução da automação extensiva em processos
pode vir a ter as mesmas consequências para os engenheiros de software. Se
sua resposta for não, justifique. Se você acha que sim, que vai reduzir as
oportunidades de emprego, é ética a resistência passiva ou ativa, pelos
engenheiros afetados, a introdução dessa tecnologia?
\end{questao}

\begin{resposta}

\textbf{A automação do processo de software e o impacto sobre o emprego:}

A resposta mais honesta é \textbf{sim, parcialmente}. A automação extensiva ---
incluindo ferramentas de geração de código por IA, testes automatizados,
integração contínua e plataformas \textit{low-code/no-code} --- tende a
deslocar uma parte significativa de tarefas de engenharia de software,
especialmente as mais rotineiras:

\begin{itemize}[nosep]
  \item Codificação de lógica de negócio simples e bem especificada;
  \item Geração de testes unitários e de integração;
  \item Refatoração e modernização de código legado;
  \item Desenvolvimento de interfaces padrão (CRUD, dashboards).
\end{itemize}

No entanto, a automação \textbf{não elimina} a necessidade de engenheiros para:
\begin{itemize}[nosep]
  \item Elicitar e interpretar requisitos complexos e ambíguos dos stakeholders;
  \item Tomar decisões arquiteturais com implicações de longo prazo;
  \item Garantir segurança, privacidade e conformidade regulatória;
  \item Validar criticamente a saída das ferramentas de IA (que cometem erros
        sistemáticos não óbvios);
  \item Sistemas críticos onde verificação formal é obrigatória.
\end{itemize}

O efeito mais provável é uma \textbf{transformação da profissão}, similar ao que
ocorreu com outras disciplinas de engenharia com a chegada de ferramentas CAD:
redução de papéis operacionais e elevação da demanda por papéis de julgamento,
supervisão e domínio especializado.

\textbf{Sobre a ética da resistência:}

A resistência --- passiva ou ativa --- pelos engenheiros afetados é compreensível
do ponto de vista humano, mas é \textbf{eticamente problemática} quando prejudica
terceiros ou o bem público:

\begin{itemize}
  \item \textbf{Resistência passiva (sabotagem sutil, atraso intencional):}
        viola os princípios de ética profissional de engenharia, que incluem
        a responsabilidade de agir no melhor interesse do público e dos
        empregadores. Atrasar a adoção de tecnologias que aumentam qualidade
        ou reduzem custos prejudica usuários e organizações.

  \item \textbf{Resistência ativa (lobby, organização sindical, pressão legal):}
        é \textbf{eticamente legítima} quando exercida dentro dos marcos legais
        e democráticos. Assim como trabalhadores industriais têm o direito de
        negociar condições de transição, engenheiros têm o direito de exigir
        políticas de requalificação, aviso prévio e reconhecimento da perda
        de renda, desde que não prejudiquem a segurança dos sistemas ou os
        usuários finais.
\end{itemize}

Em suma: adaptar-se, redirecionar esforços para atividades de maior valor
agregado e participar ativamente da discussão sobre políticas de transição
é mais ético e eficaz do que resistir à mudança tecnológica per se.

\end{resposta}

\vspace{10pt}\hrule\vspace{8pt}
\begin{center}
{\small \textit{Referência:} Ian Sommerville.
\textbf{Software Engineering}, 10ª ed. Pearson, 2016. Capítulo 2:
Software Processes.}
\end{center}

\end{document}
